% ============================================================
\section{数据处理}
% 本章只做全局性的数据画像：来源、规模、变量字典、整体数据质量，以及对所有问题
% 都成立的通用处理约定（如时段编号方式）——不做面向具体某一问建模需求的清洗/
% 特征决策，那些放在各问求解章各自的“N.1 数据预处理”小节中。5.1 只说明数据
% 本身长什么样，5.2 说明为了让全部问题能在统一口径下使用这批数据而做的通用处理。

\subsection{数据说明}
本题数据由题目附件1至附件4提供，均为微网运行相关的时间序列数据，采样粒度统一为10分钟（每天144个时段）。附件1给出某一典型日全天的电价、小区负载与光伏预测功率，共\textbf{144条记录}；附件2给出2025年全年每天的小区负载与光伏发电实际功率，共\textbf{52{,}560条记录}；附件3给出2025年全年每天0:00、6:00、12:00、18:00共四个时刻发布的未来24小时整点光伏发电功率预报，累计\textbf{35{,}040条}原始整点预报记录；附件4给出2025年全年每天不同时间的实际电价，同样为\textbf{52{,}560条记录}。

\subsection{数据处理}
为便于各问在统一口径下调用同一批数据，本文对原始数据做了时段编号处理。将每天划分为$t=1,2,\ldots,144$共144个时段，第$t$个时段对应的时间区间为
\begin{equation}
I_t=\big[(t-1)\Delta t,\ t\Delta t\big),\qquad \Delta t=\tfrac16\ \text{小时（10分钟）}.
\end{equation}
即时段$t$覆盖当天第$(t-1)\times10$分钟至第$t\times10$分钟。全年数据按来源日期与$t=1,\ldots,144$关联，某日24:00的采样仍归属当天（不归入次日），从而保证同一天的144个时段构成一个完整、无重叠的日历日。